% Part: second-order-logic
% Chapter: metatheory
% Section: introduction

\documentclass[../../../include/open-logic-section]{subfiles}

\begin{document}

\olfileid[fr]{sol}{met}{int}

\olsection{Introduction}

La logique du premier ordre possède plusieurs
propriétés remarquables. Nous savons qu'elle
n'est pas décidable, mais elle est du moins
axiomatisable. Il existe en effet des systèmes
de preuve corrects et complets pour cette
logique : ils définissent une relation de
!!{derivability}~$\Proves$ telle que, pour
tout ensemble d'!!{sentence}s~$\Gamma$ et
tout !!{sentence}~$!A$,
$\Gamma\Entails !A$ si et seulement si
$\Gamma\Proves !A$.
En particulier, pour un langage fixé dont
les symboles sont effectivement codés, les
!!{sentence}s valides du premier ordre sont
!!{computably enumerable}s. Il existe une
fonction calculable
$f\colon\Nat\to\Sent[L]$ dont les valeurs
sont exactement les !!{sentence}s valides de
$\Lang{L}$. En effet, on peut énumérer les
!!{derivation}s sans hypothèses, puis associer à chacune de
celles qui concluent un seul !!{sentence}
cet !!{sentence}. La logique du second ordre
est plus expressive que celle du premier ordre;
ses validités sont donc en général plus difficiles
à caractériser. Nous montrerons même que, sous
la sémantique standard, elles ne sont pas seulement
indécidables : elles ne sont pas
!!{computably enumerable}s. Il n'existe donc
aucun système de preuve effectif à la fois correct et
complet pour cette logique, bien que des systèmes
corrects mais incomplets existent et constituent
d'importants objets de recherche.

La logique du premier ordre possède encore deux
propriétés : la compacité (si chaque sous-ensemble
fini d'un ensemble~$\Gamma$ d'!!{sentence}s est
satisfaisable, alors $\Gamma$ lui-même est
satisfaisable) et, pour les langages au plus
dénombrables, la propriété de Löwenheim--Skolem
(si $\Gamma$ possède un modèle infini, elle
possède un modèle !!{denumerable}). Ces deux
propriétés échouent au second ordre, sous la
sémantique standard. Là encore, la raison est
que cette logique peut exprimer sur la taille
des !!{domain}s des faits que le premier ordre
ne peut pas exprimer.

\begin{editorial}
La source énonce l'énumération calculable des
validités et la forme dénombrable du théorème
de Löwenheim--Skolem sans préciser les conditions
sur le langage. Nous indiquons respectivement
le codage effectif des symboles et la borne
dénombrable sur le langage. Les échecs annoncés
pour la logique du second ordre concernent sa
sémantique standard.
\end{editorial}

\end{document}
